\section{四条研究主线如何连接}

\begin{table}[H]
\centering
\small
\begin{tabular}{>{\raggedright\arraybackslash}p{3cm}>{\raggedright\arraybackslash}p{3.4cm}>{\raggedright\arraybackslash}p{3.7cm}>{\raggedright\arraybackslash}p{3.4cm}}
\toprule
\textbf{方向} & \textbf{突破的瓶颈} & \textbf{核心机制} & \textbf{新的风险} \\
\midrule
Multi-Agent FT & reasoning data 同质化 & 多 generator、critic、debate & correlated error、成本 \\
DeepSeekMath-V2 & verifier 不可信 & verifier + meta-verifier & judge 共谋、域外失效 \\
Absolute Zero & 训练任务静态且昂贵 & proposer--solver 动态课程 & 无效任务、环境投机 \\
Intelligence/Watt & 能耗与云端供给限制 & workload-aware 能效与混合路由 & 质量退化、测量偏差 \\
\bottomrule
\end{tabular}
\end{table}

\begin{importantbox}{未来自我改进是四层共设计}
数据层生成新任务与轨迹；验证层阻止错误进入训练；策略层学习搜索、工具与协作；系统层提供可负担的计算。只优化其中一层，收益会被另一层瓶颈截断。
\end{importantbox}

\subsection{开放问题}

\begin{itemize}
    \item 如何度量真正的 reasoning diversity，而不是表面 embedding 差异？
    \item Meta-verifier 如何在没有专家标签的开放领域校准？
    \item 自生成任务如何保证与现实目标相关，而非形成封闭“自娱自乐”的课程？
    \item 如何让慢仿真与真实实验进入高频 RL loop，而不被 surrogate reward 劫持？
    \item Local/cloud router 如何联合考虑质量、隐私、延迟、能耗和数据驻留？
    \item 哪些经验应该写 memory，哪些应该触发 parameter update？
\end{itemize}

\subsection{本章小结}

四条主线分别扩展数据、验证、任务与计算效率，但最终必须合并为一个可审计、可持续、资源感知的 self-improvement stack。

